\documentclass{beamer}

\usepackage[russian]{babel}
\usepackage[utf8]{inputenc}
\usepackage{cmap}
\usepackage{graphicx}
\usepackage{xspace}
\usepackage{psfrag}
\usepackage{hyperref}
\usepackage[normalem]{ulem}

\beamertemplatenavigationsymbolsempty
\setbeamertemplate{footline}[frame number]
\usecolortheme{seahorse}

%%Исправляем греческие буквы в формулах
\DeclareSymbolFont{greek}{U}{eur}{m}{n}                                                        
\SetSymbolFont{greek}{bold}{U}{eur}{b}{n}                                                      
\DeclareSymbolFontAlphabet{\gr}{greek}                                                         
\DeclareMathSymbol{\alpha}{\mathord}{greek}{"0B}                                               
\DeclareMathSymbol{\beta}{\mathord}{greek}{"0C}                                                
\DeclareMathSymbol{\gamma}{\mathord}{greek}{"0D}                                               
\DeclareMathSymbol{\delta}{\mathord}{greek}{"0E}                                               
\DeclareMathSymbol{\epsilon}{\mathord}{greek}{"0F}                                             
\DeclareMathSymbol{\zeta}{\mathord}{greek}{"10}                                                
\DeclareMathSymbol{\eta}{\mathord}{greek}{"11}                                                 
\DeclareMathSymbol{\theta}{\mathord}{greek}{"12}                                               
\DeclareMathSymbol{\iota}{\mathord}{greek}{"13}                                                
\DeclareMathSymbol{\kappa}{\mathord}{greek}{"14}                                               
\DeclareMathSymbol{\lambda}{\mathord}{greek}{"15}                                              
\DeclareMathSymbol{\mu}{\mathord}{greek}{"16}                                                  
\DeclareMathSymbol{\nu}{\mathord}{greek}{"17}                                                  
\DeclareMathSymbol{\xi}{\mathord}{greek}{"18}                                                  
\DeclareMathSymbol{\pi}{\mathord}{greek}{"19}                                                  
\DeclareMathSymbol{\rho}{\mathord}{greek}{"1A}                                                 
\DeclareMathSymbol{\sigma}{\mathord}{greek}{"1B}                                               
\DeclareMathSymbol{\tau}{\mathord}{greek}{"1C}                                                 
\DeclareMathSymbol{\upsilon}{\mathord}{greek}{"1D}                                             
\DeclareMathSymbol{\phi}{\mathord}{greek}{"1E}                                                 
\DeclareMathSymbol{\chi}{\mathord}{greek}{"1F}                                                 
\DeclareMathSymbol{\psi}{\mathord}{greek}{"20}                                                 
\DeclareMathSymbol{\omega}{\mathord}{greek}{"21}                                               
\DeclareMathSymbol{\varepsilon}{\mathord}{greek}{"22}                                          
\DeclareMathSymbol{\vartheta}{\mathord}{greek}{"23}                                            
\DeclareMathSymbol{\varpi}{\mathord}{greek}{"24}                                               
\let\varrho\rho                                                                                
\let\varsigma\sigma                                                                            
\DeclareMathSymbol{\varphi}{\mathord}{greek}{"27}

% === ТИТУЛЬНЫЙ ЛИСТ ===
\title{Разработка гибридной рекомендательной системы на основе аудио-эмбеддингов с использованием механизма Two-Tower Attention}
\subtitle{Защита курсового проекта}
\author{Малютин~А.~А.}
\institute{Институт прикладной математики и компьютерных наук ТГУ}
\date{23 июня 2026 г.}

\begin{document}

\begin{frame}
\maketitle
\end{frame}

% === СЛАЙД 1 ===
\begin{frame}
\frametitle{Актуальность и цель работы}
\begin{itemize}
    \item \textbf{Актуальность:} Пользователи сталкиваются с миллионами треков, что вызывает необходимость учитывать индивидуальные музыкальные предпочтения[cite: 2].
    \item \textbf{Проблематика:} Классические системы уязвимы к проблеме <<холодного старта>> при добавлении новых композиций.
    \item \textbf{Цель работы:} Разработать рекомендательную систему на основе аудио-эмбеддингов с применением архитектуры Two-Tower Attention для решения проблемы холодного старта.
\end{itemize}
\end{frame}

% === СЛАЙД 2 ===
\begin{frame}
\frametitle{Анализ существующих решений}
\begin{itemize}
    \item \textbf{Коллаборативная фильтрация:} Принцип <<Люди со схожими вкусами любят схожую музыку>>[cite: 2].
    \begin{itemize}
        \item Плюсы: Высокая точность[cite: 2]. Минусы: <<Холодный старт>> для новых треков и пользователей[cite: 2].
    \end{itemize}
    \item \textbf{Контентные методы:} Принцип <<Если вам нравится трек, вам понравятся похожие на него>>[cite: 2].
    \begin{itemize}
        \item Плюсы: Нет проблемы холодного старта для новых треков[cite: 2]. Минусы: Не учитывает вкусы сообщества[cite: 2].
    \end{itemize}
    \item \textbf{Гибридные методы:} Объединение лучших качеств обоих методов[cite: 2].
    \begin{itemize}
        \item Плюсы: Максимальная точность, решает проблему холодного старта[cite: 2].
    \end{itemize}
\end{itemize}
\end{frame}

% === СЛАЙД 3 ===
\begin{frame}
\frametitle{Объект исследования: Датасет Yandex Yambda}
Yambda --- это крупнейший открытый датасет музыкальных взаимодействий от Яндекс.Музыки[cite: 2].

\vspace{0.3cm}
\textbf{Ключевые характеристики:}
\begin{itemize}
    \item Объем: 4.79 млрд взаимодействий (прослушивания, лайки, дизлайки)[cite: 2].
    \item Охват: 1 млн пользователей и 9.39 млн треков[cite: 2].
    \item Признаки: Аудио-эмбеддинги --- математическое описание звучания треков[cite: 2].
    \item Разметка: Органические vs рекомендательные прослушивания[cite: 2].
\end{itemize}
\end{frame}

% === СЛАЙД 4 ===
\begin{frame}
\frametitle{Архитектура модели: Two-Tower Attention}
Предложенная нейросеть разделяет обработку профилей и контента:
\begin{itemize}
    \item \textbf{User Tower (Башня пользователя):} Векторизация истории профиля слушателя через слой \texttt{Embedding} и нормализацию (\texttt{LayerNorm}).
    \item \textbf{Item Tower (Башня трека):} Обработка плоских аудио-эмбеддингов (размерность 128) через полносвязные слои.
    \item \textbf{Latent Sequence Attention:} Инновационный блок модели. Аудио-вектор искусственно переформируется в последовательность токенов. \textit{Multi-Head Attention} (4 головы) позволяет сети фокусироваться на независимых характеристиках звука (ритм, тембр, вокал).
    \item \textbf{Обучение:} Контрастивный подход с In-Batch Negatives.
\end{itemize}
\end{frame}

% === СЛАЙД 5 ===
\begin{frame}
\frametitle{Реализация и веб-интерфейс}
\begin{itemize}
    \item \textbf{Сценарий:} Рекомендация треков для пользователя с историей прослушиваний[cite: 2].
    \item \textbf{Процесс ранжирования:} 
    \begin{enumerate}
        \item Извлечение профиля пользователя.
        \item Анализ музыкальных признаков.
        \item Балансировка: итоговый скор вычисляется как сумма предсказания нейросети и глобальной популярности трека ($NN_{scaled} + \alpha \cdot Pop_{scaled}$).
    \end{enumerate}
    \item \textbf{Интерфейс (Streamlit):} Для интерактивной демонстрации реализован веб-стенд. Применены алгоритмы кэширования матриц на диск, что позволяет визуализировать смену приоритетов выдачи в реальном времени без задержек.
\end{itemize}
\end{frame}

% === СЛАЙД 6 ===
\begin{frame}
\frametitle{Заключение и источники}
\textbf{Основные результаты:}
\begin{itemize}
    \item Спроектирована архитектура Two-Tower Attention для глубокого анализа аудио-эмбеддингов.
    \item Эффективно решена проблема <<холодного старта>> за счет гибридной балансировки.
    \item Разработан рабочий прототип системы для демонстрации алгоритмов.
\end{itemize}

\vspace{0.3cm}
\textbf{Список основных источников:}
\begin{enumerate}
    \item Covington P. et al. Deep Neural Networks for YouTube Recommendations // RecSys '16.
    \item Vaswani A. et al. Attention Is All You Need // NIPS '17.
    \item Yandex Research. Yambda Dataset // GitHub, 2026.
\end{enumerate}
\end{frame}

\end{document}